\pdfoutput=1

\usefonttheme{professionalfonts}  % get CMS font in beamer
\usepackage[scaled=0.87]{helvet}    %helvetica
% \usetheme[darktitle]{UAFrederik}
\usetheme{cern}

% Make alignment in itemize justified again
\usepackage{ragged2e}
\makeatletter
\renewcommand{\itemize}[1][]{%
  \beamer@ifempty{#1}{}{\def\beamer@defaultospec{#1}}%
  \ifnum \@itemdepth >2\relax\@toodeep\else
    \advance\@itemdepth\@ne
    \beamer@computepref\@itemdepth% sets \beameritemnestingprefix
    \usebeamerfont{itemize/enumerate \beameritemnestingprefix body}%
    \usebeamercolor[fg]{itemize/enumerate \beameritemnestingprefix body}%
    \usebeamertemplate{itemize/enumerate \beameritemnestingprefix body begin}%
    \list
      {\usebeamertemplate{itemize \beameritemnestingprefix item}}
      {\def\makelabel##1{%
          {%
            \hss\llap{{%
                \usebeamerfont*{itemize \beameritemnestingprefix item}%
                \usebeamercolor[fg]{itemize \beameritemnestingprefix item}##1}}%
          }%
        }%
      }
  \fi%
  \beamer@cramped%
  \justifying% NEW
  %\raggedright% ORIGINAL
  \beamer@firstlineitemizeunskip%
}
\makeatother



%\mode<handout>{
% \usepackage{pgfpages}
% \pgfpagesuselayout{4 on 1}[a4paper, border shrink=5mm, landscape]
%}



% key - value options
\pgfkeys{
 	/marker/.is family, /marker,
	default/.style = {color=LimeGreen,inner sep=1.5mm},
	default2/.style = {color=LimeGreen,inner sep=0.5pt},
	color/.estore in = \markercolour,
	inner sep/.estore in = \markerinnersep
}

% put this before the first block on a slide to gain some space
\newcommand{\firstblock}{\vspace*{-2em}}

% \highlight<>[colour][text]  puts text in colour
\newcommand<>{\highlight}[2][uared]{{\color#3{#1}#2}}
% \markermath<>[options]{formula}{nodename} puts a box around the formula. Nodename specifies an (unique!)
% name for the node, such that it can later be connected with an arrow. Options are color=... and inner sep=... 
\newcommand<>{\markermath}[3][]{
	\pgfkeys{/marker, default, #1}%
	\alt#4
	{\tikz[baseline,remember picture]
		\node[rectangle, fill=\markercolour,rounded corners, inner sep=\markerinnersep,anchor=base] (#3)  {$\displaystyle#2$};}
	{\tikz[baseline,remember picture]
		\node[inner sep=\markerinnersep,anchor=base]  (#3)  {$\displaystyle#2$};}
}
% \arrowmarkerstart{text}{nodename}  is an invisible node which is the start for an arrow
\newcommand{\arrowmarkerstart}[3][]{
	\pgfkeys{/marker, default2, #1}%
	\tikz[baseline,remember picture]
		\node[rectangle,rounded corners, inner sep=\markerinnersep,anchor=base] (#3) {#2};
}
% \arrowmarker[bend options]{node1}{node2} is an arrow from node1 to node2
\newcommand<>{\arrowmarker}[3][bend left]{
	\tikz[remember picture, overlay] \path[->]#4 (#2) edge [#1] (#3);
}

% \beamerdefaultoverlayspecification{<+->}
\setbeamercovered{transparent}



